\documentclass[conference,a4paper]{IEEEtran}
\usepackage[utf8]{inputenc}
\usepackage[T1]{fontenc}
\usepackage[brazil]{babel}
\usepackage{array,amsmath,amssymb,booktabs,tabularx,graphicx,tikz}
\usetikzlibrary{arrows.meta,positioning}
\usepackage[hidelinks]{hyperref}
\usepackage{xurl}
\newcolumntype{Y}{>{\raggedright\arraybackslash}X}
\newcolumntype{P}[1]{>{\raggedright\arraybackslash}p{#1}}
\newcommand{\campo}[1]{\textbf{#1.}}
\newcommand{\ctt}[1]{\texttt{#1}}
\addto\captionsbrazil{\renewcommand{\IEEEkeywordsname}{Index Terms}}
\setlength{\emergencystretch}{2em}
\pagestyle{plain}
\raggedbottom
\author{\IEEEauthorblockN{Lucas Cerqueira Portela e Bruno Prado dos Santos}
\IEEEauthorblockA{Engenharia de Computação -- CEFET-MG\\Campus Divinópolis, Minas Gerais, Brasil}}


\title{Implementação instrumentada de um UX-Agent para lembretes e registros de tomada de medicamentos}
\begin{document}
\maketitle
\begin{abstract}
Este trabalho apresenta uma aplicação web local para a execução de tarefas fictícias de lembrete e registro de tomada por um agente proxy baseado em modelo de linguagem. O sistema compartilha a representação pública de tela e ações entre a interface e a observação textual. A execução obedece às transições da conversa e o verificador externo avalia o desfecho. Foram verificadas a instrumentação, a ajuda sob demanda e a ausência dos campos privados conhecidos na observação, por meio de ações dirigidas nas tarefas de calibração. Na interface revisada, cinco tentativas com três modelos produziram um sucesso em C2. Em C4, o agente voltou repetidamente à ajuda ou desistiu. As observações orientam a revisão da interface e não substituem uma avaliação com pessoas. A comparação com KLM distingue contagem de ações, referência de especialista e tempo computacional.
\end{abstract}
\begin{IEEEkeywords}
UX-Agent, instrumentação, comunicabilidade, KLM, dados sintéticos.
\end{IEEEkeywords}
\section{Introdução}
Como produzir registros auditáveis de tarefas sem fornecer ao agente informação que a persona não teria? A P3 transforma um design instrumentado em aplicação e registra a cadeia entre ação, resposta pública e avaliação externa~\cite{p3}. O aplicativo não coleta dados de saúde e usa artefatos fictícios.

\section{Modelagem herdada e decisões}
A2 definiu Helena, personagem idosa que lê português, tem pouca experiência em aplicativos de medicamentos e apoio familiar, C2 (comparar ocorrências) e C4 (interpretar pendência com ajuda)~\cite{a2}. V5/V8 e suas HTAs foram preservadas na reserva. Os métodos das quatro tarefas, a hierarquia CTT, a conversa MoLIC, a matriz de rupturas e a ajuda integram E1; apenas C2/C4 foram instrumentadas como cenários executáveis para esta fase.

\begin{table}[ht]
\caption{Decisões de implementação}\small
\begin{tabularx}{\columnwidth}{@{}YYP{.19\columnwidth}@{}}\toprule
Decisão & Justificativa & Status\\\midrule
Renderizador compartilhado & A mesma fonte de signos alimenta tela e observação & Princípio de design\\
Ações com IDs opacos & Reduz metadados internos no canal de decisão & Princípio de design\\
A/B preservam controles & Comparar rótulos e persistência de feedback no mesmo desenho & Hipótese operacional\\
Cartão sob demanda & Registrar entrega, manter previsibilidade e permitir retorno & Princípio de design\\
Perfis com parâmetros externos & Representar hipóteses replicáveis e ablação técnica & Hipótese operacional\\\bottomrule
\end{tabularx}
\end{table}

\section{Arquitetura e acoplamento}
O motor lê um conjunto versionado de transições inspiradas na MoLIC~\cite{molic}\cite[pp. 9--31]{guiamolic}. O renderizador produz somente título, linhas públicas e controles. O tradutor associa cada controle a um ID opaco de uso único e fornece apenas as ações autorizadas pelo estado atual. O adaptador do modelo prepara uma chamada com uma resposta estruturada de seleção de ação. O executor verifica a combinação de ação e observação. IDs de banco, modelos de tarefas, ordenação analítica, desfecho correto e critérios não entram na requisição. A interface local fica restrita ao dispositivo e mostra a auditoria após o encerramento.

\campo{Mecanismos} M1 controla omissão de sinais compactos por visita quando desconhecidos. M2 mantém blocos já exibidos, M3 retém marcos e M4 guarda as associações do cartão solicitado. Os esquemas de referência usam $q=0{,}35$, $k=a=1$ e $l=5$. O amostrador usa grades discretas da A2 e semente por sessão. O controle neutro desliga a omissão e a retenção limitada e desativa a aprendizagem cumulativa. A ajuda permanece disponível. Entre as duas tentativas da sessão, M2/M3 reiniciam e M4 permanece. Não se introduz atraso comportamental sem sustentação empírica. Tempos registrados são da execução.

\campo{Verificador} Usa estados dos relatos, consulta às ocorrências, resposta selecionada e entrega de ajuda quando exigida. Desistência é insucesso; falha técnica que impede decidir é desconhecida. O encerramento acontece mesmo com resposta incompleta. O sucesso não usa afirmação textual do agente.

\section{Instrumentação e critérios}
Cada evento tem sessão, sequência, tempo UTC, tipo e origem; quando aplicável, tarefa, variante, posição, cena, fala, resposta e dados técnicos. A ajuda registra pedido, entrega, camada e retorno. Os padrões operacionais MAC incluem pedir ajuda, cancelar a confirmação, repetir navegação, desistir e concluir sem sucesso; são candidatos à inspeção~\cite[pp. 31--38]{mac}. O pacote E3 inclui a trilha e o dicionário. O E3 guarda os prompts e observações fornecidos pela aplicação à CLI, as respostas disponíveis e as latências, sem credenciais; o contexto interno gerenciado pelo serviço não é exportado.

\section{Verificação e piloto}
A verificação técnica executou quatro sessões de duas tentativas com ações dirigidas, cobrindo C2/C4, A/B e as duas ordens. Oito verificações automatizadas exercitaram rota por U/D, erro recuperado, obrigatório de ajuda, reinício de filas, aprendizagem, rejeição e falha, cancelamento e reserva. Nas oito tentativas dirigidas, houve oito fechamentos com sucesso externo. Esses caminhos foram escolhidos para conferir a determinística do instrumento e não estimam taxa de sucesso de nenhum público.

\begin{table}[ht]
\caption{Trajetórias dirigidas e referência ideal}\small
\begin{tabularx}{\columnwidth}{@{}lYYYY@{}}\toprule
Tarefa & Sucesso & Erros / op. & Ações / tentativa & Ajuda / tentativa\\\midrule
C2 pela rota H/H & 4/4 & 0/4 & 6 & 0\\
C4 & 4/4 & 0/4 & 4 & 1\\\bottomrule
\end{tabularx}
\end{table}

As contagens são compatíveis com H/H de C2 (6 ações, 10,50~s) e C4 (4 ações, 6,55~s), referências KLM de especialista com mouse~\cite[pp. 397--400]{card}. U/H, rota alternativa de C2, usa 5 ações e 9,20~s na referência. O KLM não foi comparado a latências do LLM. 
\begin{table}[ht]\caption{Piloto com LLM: configuração e tarefa}\small
\begin{tabularx}{\columnwidth}{@{}lYYYY@{}}\toprule
Condição & S/N & Desc. & A méd. & H méd.\\\midrule
Fixa/C2 & 7/7 & 0 & 9.43 & 0.00\\
Fixa/C4 & 0/8 & 1 & 39.00 & 19.62\\
Neutra/C2 & 1/1 & 0 & 8.00 & 0.00\\
Neutra/C4 & 0/1 & 0 & 40.00 & 24.00\\
\bottomrule\end{tabularx}\end{table}
\subsection{Escopo e resultados do piloto}
A coleta com LLM foi encerrada por disponibilidade de franquia e prazo: a meta anterior de 16 sessões foi cancelada, sem alteração de parâmetros ou tarefas. O enunciado exige o registro-piloto, sem fixar quantidade mínima. Foram preservadas nove sessões (oito na referência fixa e uma neutra), com 17 tentativas avaliadas: oito sucessos, oito insucessos por limite e uma tentativa desconhecida por falha técnica. Uma tarefa posterior não iniciada não integra o denominador. A análise abrange somente os arquivos persistidos; trabalho interrompido antes da gravação não foi reconstruído. Smoke e trajetórias dirigidas não foram agregados à amostra.

$S$ representa sucessos confirmados por tentativas iniciadas, com desconhecidos informados à parte. $A$ e $H$ são médias de ações executadas e entregas de ajuda, incluindo a trajetória parcial desconhecida. Nas oito tentativas C2, $R_e=0/8$ oportunidades; em C4, o denominador funcional foi zero, portanto $R_e$ é indefinido, não zero. A ausência de erros E1--E3 não implica sucesso: a tarefa pode terminar sem fornecer a resposta.

As oito tentativas C2 tiveram sucesso externo. As oito tentativas C4 conhecidas atingiram 40 ações sem conclusão satisfatória; a nona foi interrompida tecnicamente após 32 ações. Houve repetição de pedidos de ajuda. A configuração neutra também apresentou esse padrão; não há base para atribuí-lo aos mecanismos da persona. Cada decisão ocorre sem histórico conversacional automático; o contexto entre passos depende do estado externo. A retenção da ajuda e a política de escolha requerem investigação posterior, sem ajustes neste conjunto.

Em C2, as médias de ações excederam a referência ideal: 9,43 na condição fixa e 8,00 na neutra, frente a 5 ou 6 no KLM. Em C4, 39,00 e 40,00 ações não são tempos nem trajetórias equivalentes à referência bem-sucedida de 4 ações. O KLM de especialista fornece 9,20--10,50~s em C2 e 6,55~s em C4; tempos do serviço e da CLI são custo computacional, não tempo humano modelado.

A mediana das latências de decisões recebidas foi 3.95~s; os arquivos contêm 427 requisições persistidas. O adaptador Codex usou login ChatGPT, saída estruturada e modelo padrão da conta. O identificador efetivo do modelo não foi emitido. Os parâmetros do adaptador alternativo da Responses API não foram usados nesta coleta.
A amostra é pequena, desigual e dependente dentro das sessões. Não foram realizados testes de significância ou alegações de superioridade entre condições. O piloto demonstra execução auditável, não validação da persona. Os prompts, observações da aplicação, seleções e latências constam do E3. A CLI adiciona contexto gerenciado não exportado pelo adaptador; o pacote não captura esse contexto interno integral e não permite controlar completamente o serviço. O vídeo E5 reproduz uma trajetória real de C2 registrada, identificada como reprodução, sem novas chamadas ao modelo.


\section{Ameaças à validade e conclusão}
A transcrição textual preserva o conteúdo sem replicar leitura visual e movimentação. A máscara de sinais em A introduz uma diferença modelada. Os mecanismos vêm de hipóteses da A2 e os dados disponíveis não fornecem comparação humana diretamente compatível. Contagens simuladas, respostas estruturadas do modelo e referência KLM não estimam adesão, ingestão, satisfação ou experiência da população. Foi possível verificar localmente a cadeia observação--ação--evento--desfecho. As execuções demonstram a cadeia de registro. Os modelos e limites diferem entre as rodadas, e o contexto interno da CLI não pode ser reconstituído integralmente. A repetição da ajuda em C4 é um ponto para investigação, e o piloto não valida a persona.

\bibliographystyle{IEEEtran}
\bibliography{referencias}
\end{document}

